\chapter{Project 1: A Transformer Language Model from Scratch}\label{ch:proj1}

Part~VIII closes the book with five capstone projects. Each project chapter is a specification
and a guide: what to build, how to set it up in this repository, how your work is checked, which
chapters teach the ideas you need, where the traps are, and how to design, log and report the
experiments. Project~1 brings together Parts~II and~III. You build every piece of a modern
language-model training stack yourself, then use it to answer questions empirically.

\begin{keyidea}[About the project guides]
The project guides deliberately stop short of solutions, so that you build everything yourself.
They specify interfaces and acceptance tests, explain the background, and suggest how to debug
and how to design experiments, but they do not contain the implementation. The book's
\code{llmsys} package does implement many of the same concepts (a BPE tokenizer, attention, AdamW,
a training loop): it is the \emph{reference} that the acceptance tests compare your code against.
Do not import or copy it. Reading it after you have a working version of your own is a good way
to compare designs.
\end{keyidea}

% =====================================================================
\section{What you build}\label{sec:proj1:overview}

Starting from tensors and a handful of container classes, you will build:

\begin{enumerate}
  \item a \textbf{byte-level BPE tokenizer}: a trainer that learns merges from a corpus, and an
        encoder/decoder with special tokens and a streaming mode for files larger than memory;
  \item a \textbf{pre-norm Transformer LM} with RMSNorm, SwiGLU feed-forward layers, rotary
        position embeddings and causal multi-head attention, built from your own linear and
        embedding layers;
  \item the \textbf{training system}: a numerically stable cross-entropy, AdamW, a warmup + cosine
        schedule, global-norm clipping, a random-window sampler over memory-mapped token files,
        checkpoint save/restore, a configurable training script and experiment logging;
  \item \textbf{decoding} with temperature and nucleus (top-$p$) sampling;
  \item a set of \textbf{experiments} on TinyStories (Eldan and Li, 2023) and a sample of
        OpenWebText (Gokaslan and Cohen, 2019), written up as a short report.
\end{enumerate}

\paragraph{The from-scratch rule.} From \code{torch.nn}, \code{torch.nn.functional} and
\code{torch.optim} you may use only \code{nn.Parameter}, the container modules (\code{Module},
\code{ModuleList}, \code{ParameterDict}, \ldots) and the \code{Optimizer} base class. Tensor
operations, \code{einsum}/\code{einops}, and \code{torch.nn.init} are fine. Once your own version
passes its acceptance test you may switch individual pieces to fused library kernels for speed,
for instance in the speed-run challenge (\cref{sec:proj1:speedrun}); say so in your report.

\paragraph{Deliverables.} (1) Your code, passing the acceptance tests of \cref{sec:proj1:setup};
(2) trained checkpoints and the run directories of every experiment; (3) a report of 6--10 pages
answering the questions in \cref{sec:proj1:experiments}, with an experiment log.

\begin{figure}[t]
\centering
\resizebox{\linewidth}{!}{%
\begin{tikzpicture}[
  comp/.style={draw, rounded corners=3pt, minimum height=0.9cm, align=center, font=\small, line width=0.6pt, text width=3.3cm},
  ref/.style={font=\scriptsize, text=srcgray, align=center},
  arr/.style={-{Stealth[length=2mm]}, line width=0.6pt}]
\node[comp, fill=fdata!18, draw=keygreen] (tok) at (0,0) {M1: BPE trainer +\\tokenizer};
\node[ref, below=1pt of tok] {Ch.~3, Lab~3; \S\ref{sec:proj1:bpe}};
\node[comp, fill=fdata!18, draw=keygreen] (ids) at (0,-2.9) {E1: uint16 token\\files};
\node[ref, below=1pt of ids] {Ch.~3, Ch.~9};
\node[comp, fill=fcompute!18, draw=sysblue] (mods) at (4.6,0) {M2: linear, embedding,\\RMSNorm, SwiGLU};
\node[ref, below=1pt of mods] {Ch.~6, Ch.~7};
\node[comp, fill=fcompute!18, draw=sysblue] (attn) at (4.6,-2.9) {M2: RoPE, softmax,\\causal MHA};
\node[ref, below=1pt of attn] {Ch.~4, Ch.~6};
\node[comp, fill=fcompute!18, draw=sysblue] (lm) at (9.2,-1.45) {M2: pre-norm block\\$\to$ Transformer LM};
\node[ref, below=1pt of lm] {Ch.~5, Ch.~6};
\node[comp, fill=fmemory!20, draw=napkin] (opt) at (13.8,0) {M3: cross-entropy, AdamW,\\schedule, clipping};
\node[ref, below=1pt of opt] {Ch.~9, Ch.~10};
\node[comp, fill=fmemory!20, draw=napkin] (loop) at (13.8,-2.9) {M4: sampler, checkpoints,\\training script};
\node[ref, below=1pt of loop] {Ch.~9, Lab~12};
\node[comp, fill=fcomm!16, draw=labpurple] (exp) at (18.4,-1.45) {E2--E7: sweeps,\\ablations, OWT,\\speed run};
\node[ref, below=1pt of exp] {Ch.~7, Ch.~9, \S\ref{sec:proj1:experiments}};
\draw[arr] (tok)--(ids); \draw[arr] (mods)--(lm); \draw[arr] (attn)--(lm);
\draw[arr] (lm)--(opt); \draw[arr] (lm)--(loop);
\draw[arr, black!60] ([yshift=-0.42cm]ids.south) -- ++(0,-0.35) -| node[pos=0.25,below,font=\scriptsize]{token ids feed the sampler} ([yshift=-0.42cm]loop.south);
\draw[arr] (opt)--(exp); \draw[arr] (loop)--(exp);
\end{tikzpicture}}
\caption{The components of Project~1 in dependency order, labelled with milestones (M) and
experiments (E), and the chapters that teach the underlying ideas.}
\label{fig:proj1:components}
\end{figure}

% =====================================================================
\section{Setting up}\label{sec:proj1:setup}

\paragraph{Your workspace.} Create a package of your own anywhere (a separate git repository
is best) and copy the interface file into it:
\begin{lstlisting}[style=shell]
mkdir -p ~/p1 && cd ~/p1 && git init
python -m venv .venv && source .venv/bin/activate
pip install -r <book>/code/requirements.txt          # torch, numpy, regex, pytest, ...
cp <book>/code/projects/p1_lm/p1_template.py p1_impl.py
python <book>/code/projects/p1_lm/check_env.py      # devices, bf16/TF32, RAM, disk
\end{lstlisting}
\texttt{p1\_impl.py} lists every function and class the tests need, with docstrings that fix the
conventions (\cref{tab:proj1:spec}). Every body raises \code{NotImplementedError}. Keep your
real code in modules of your own and make \texttt{p1\_impl.py} a thin facade that imports from
them, so your design is not constrained by the interface.

\paragraph{Acceptance tests.} The tests in \texttt{code/projects/p1\_lm/test\_p1\_acceptance.py}
load the file named by \texttt{P1\_IMPL} and check it (\cref{fig:proj1:harness}). Tokenizer tests
use property checks: structure, first merge, round trips, special tokens, streaming. Numerical
components are compared with PyTorch and with the \code{llmsys} reference at tight tolerances. The
model is tested for shape, causality, initial loss and the ability to overfit one batch. Without
\texttt{P1\_IMPL} the tests are skipped.
\begin{lstlisting}[style=shell]
P1_IMPL=~/p1/p1_impl.py python -m pytest <book>/code/projects/p1_lm -q
P1_IMPL=~/p1/p1_impl.py python -m pytest <book>/code/projects/p1_lm -q -k rope   # one test
\end{lstlisting}

\begin{figure}[t]
\centering
\resizebox{\linewidth}{!}{%
\begin{tikzpicture}[
  b/.style={draw, rounded corners=2pt, minimum height=0.85cm, align=center, font=\small, text width=4.2cm, line width=0.6pt},
  arr/.style={-{Stealth[length=2mm]}, line width=0.6pt}]
\node[b, fill=black!5] (t) at (0,0) {\texttt{test\_p1\_acceptance.py}};
\node[b, fill=fmemory!20, draw=napkin] (a) at (5.6,0) {\texttt{p1\_impl.py}\\thin facade};
\node[b, fill=fcompute!18, draw=sysblue] (p) at (11.2,0) {your modules\\(tokenizer, model, train)};
\node[b, fill=black!5] (r) at (0,-1.9) {references: PyTorch,\\\code{llmsys}, properties};
\draw[arr] (t) -- node[above,font=\scriptsize]{\texttt{P1\_IMPL}} (a);
\draw[arr] (a) -- node[above,font=\scriptsize]{imports} (p);
\draw[arr, dashed] (r) -- node[left,font=\scriptsize]{expected values} (t);
\draw[arr, dashed] (p.south) |- node[pos=0.72,below,font=\scriptsize]{outputs compared} (r.east);
\end{tikzpicture}}
\caption{The acceptance harness. The tests reach your code only through the interface file, so
you are free to organise the real implementation as you like.}
\label{fig:proj1:harness}
\end{figure}

\paragraph{Data.} TinyStories is distributed on the Hugging Face Hub
(dataset \ident{roneneldan/TinyStories}; use the \ident{TinyStoriesV2-GPT4} train and
validation text files, whose stories are separated by \texttt{<|endoftext|>}). For web text,
take a sample of a few gigabytes of OpenWebText, for example from the
\ident{Skylion007/openwebtext} dataset on the Hub, and write it as
one text file with \texttt{<|endoftext|>} between documents, and hold out a validation file. Budget
tens of gigabytes of disk for text plus the \code{uint16} token arrays.

\paragraph{Compute.} Everything up to E5 runs on a laptop with the reduced budgets of
\cref{sec:proj1:outside}; a single modern GPU runs the full budgets in minutes per run.
\code{budget.py} converts between tokens, steps, batch size and context length, and estimates
wall-clock time from your measured throughput.

% =====================================================================
\section{Specification and milestones}\label{sec:proj1:map}

\Cref{tab:proj1:spec} is the interface, in the order in which you should build it. ``Teaches''
points to where the book develops the idea; the acceptance tests define ``done''.

{\footnotesize
\begin{longtable}{@{}>{\raggedright\arraybackslash}p{0.07\linewidth}>{\raggedright\arraybackslash}p{0.22\linewidth}>{\raggedright\arraybackslash}p{0.36\linewidth}>{\raggedright\arraybackslash}p{0.27\linewidth}@{}}
\caption{Project~1 interface and milestones.}\label{tab:proj1:spec}\\
\toprule
\textbf{MS} & \textbf{Interface} & \textbf{Contract} & \textbf{Teaches}\\
\midrule\endfirsthead
\toprule
\textbf{MS} & \textbf{Interface} & \textbf{Contract} & \textbf{Teaches}\\
\midrule\endhead
M1 & \code{train\_bpe(text, vocab\_size, specials)} & ids 0--255 are bytes, then merges in creation order, specials last; specials split text and are not counted; ties go to the greater pair; deterministic & \cref{sec:text:bpe}, Lab~3; \cref{sec:proj1:unicode,sec:proj1:bpe}\\
M1 & \code{Tokenizer(vocab, merges, specials)}: \code{encode}, \code{encode\_\allowbreak iterable}, \code{decode} & exact round trips (Unicode, whitespace, specials); overlapping specials resolved longest-first; streaming equals \code{encode}; invalid UTF-8 decodes to U+FFFD & \cref{sec:text:special,sec:text:wild}; \cref{sec:proj1:stream}\\
M2 & \code{softmax}, \code{cross\_entropy} & stable for logits near $10^4$; mean NLL over all leading dims & \cref{sec:attention:simple}, \cref{sec:optim:loss}\\
M2 & \code{rmsnorm(x, w, eps)} & fp32 statistics, output in input dtype & \cref{sec:modern:rmsnorm}\\
M2 & \code{apply\_rope(x, positions, theta)} & rotate-half pairing; scores depend only on relative position & \cref{sec:modern:rope}, \cref{sec:archdesign:pos}\\
M2 & \code{attention(q, k, v, mask)} & any leading dims; boolean mask, True = allowed & \cref{sec:attention:causal,sec:attention:mha}\\
M2 & \code{TransformerLM(...)} & pre-norm, SwiGLU, RoPE, causal, no biases; loss at init $\approx\ln V$; can overfit one batch & \cref{sec:gpt:block}, \cref{sec:modern:block,sec:modern:model}, \cref{sec:archdesign:init}\\
M3 & \code{AdamW} & matches \code{torch.optim.AdamW} trajectories & \cref{sec:optim:api,sec:optim:adam}\\
M3 & \code{lr\_schedule}, \code{clip\_grad\_norm} & 0-based warmup + cosine + floor; global norm, pre-clip norm returned & \cref{sec:optim:sched,sec:optim:clip}\\
M4 & \code{get\_batch}, \code{save/load\_\allowbreak checkpoint} & random windows from a \code{uint16} array; checkpoint restores model and optimizer so training continues identically & \cref{sec:optim:data,sec:optim:ckpt}\\
M5 & your training script & config + CLI overrides, memmapped train/val data, periodic validation, checkpoints, JSONL logs with step and wall time & \cref{sec:optim:loop,sec:optim:explog}, \cref{sec:pretrain:trainfn}\\
M6 & your decoder & temperature, top-$p$, stop at \texttt{<|endoftext|>} or a length limit & \cref{sec:pretrain:decoding}, \cref{sec:scaling:topp}\\
\bottomrule
\end{longtable}}

% =====================================================================
\section{Background: Unicode, pre-tokenisation and tokenizers at scale}\label{sec:proj1:background}

Most of the project rests on Chapters 3--13. This section adds the tokenizer details that a
production-quality M1 needs, and that \cref{ch:text} only touches on.

\subsection{Unicode, code points and UTF-8}\label{sec:proj1:unicode}

Unicode assigns each character an integer \emph{code point} (\code{ord}/\code{chr} in Python);
the standard defines roughly 160,000 characters across a space of $17\times2^{16}$ code points.
Some code points are control characters that print as nothing visible. An \emph{encoding}
turns code points into bytes. UTF-8 is variable-length:

\begin{center}\small
\begin{tabular}{@{}lll@{}}
\toprule
code point range & bytes & bit pattern\\
\midrule
U+0000--U+007F & 1 & \texttt{0xxxxxxx} (ASCII unchanged)\\
U+0080--U+07FF & 2 & \texttt{110xxxxx 10xxxxxx}\\
U+0800--U+FFFF & 3 & \texttt{1110xxxx 10xxxxxx 10xxxxxx}\\
U+10000--U+10FFFF & 4 & \texttt{11110xxx 10xxxxxx 10xxxxxx 10xxxxxx}\\
\bottomrule
\end{tabular}
\end{center}

Three consequences matter for tokenizers. First, a byte-level vocabulary needs only 256 base
symbols and can never meet an out-of-vocabulary input. Second, a single byte is not a character:
continuation bytes (\texttt{10xxxxxx}) are meaningless on their own, so any code that decodes
token-by-token or byte-by-byte must buffer incomplete sequences (\cref{sec:text:wild} builds an
incremental detokenizer for this reason). Third, not every byte sequence is valid UTF-8, so a
decoder needs an error policy. This project uses replacement with U+FFFD,
\code{bytes.decode("utf-8", errors="replace")}. Compare the same strings under UTF-8, UTF-16 and
UTF-32 to see why byte-level tokenizers are almost always built on UTF-8.

\subsection{Pre-tokenisation with the GPT-2 pattern}

The book's \code{ByteBPETokenizer} (\cref{sec:text:bpe}) uses a simplified ASCII-only pattern.
This project uses the pattern introduced with GPT-2 (Radford et al., 2019), in its common modern
spelling, which needs Unicode character classes:
\begin{lstlisting}[style=plain]
'(?:[sdmt]|ll|ve|re)| ?\p{L}+| ?\p{N}+| ?[^\s\p{L}\p{N}]+|\s+(?!\S)|\s+
\end{lstlisting}
Python's built-in \code{re} does not support \verb|\p{L}| (``any letter in any script'') or
\verb|\p{N}| (``any number''); the third-party \code{regex} module does. Read the alternatives
left to right: English contractions; an optional space followed by a run of letters; the same for
digits; the same for other non-space symbols; whitespace that is \emph{not} followed by a
non-space character (the negative lookahead \verb|\s+(?!\S)| leaves the last space of a run
attached to the next word, so ``\verb*|a   b|'' yields \verb*|a|, \verb*|  |, \verb*| b|); and
finally any remaining whitespace. Use \code{finditer} rather than \code{findall} when counting,
so you never hold a list of every pre-token of a multi-gigabyte file in memory.

\subsection{Scaling BPE training to a large file}\label{sec:proj1:bpe}

\Cref{sec:text:bpe} derives BPE, and Lab~3 measures how much faster an incremental trainer is than
the naive one. At the scale of TinyStories and an OpenWebText sample (gigabytes of text), two
system-level issues dominate.

\paragraph{Pre-tokenisation is embarrassingly parallel, if you cut the file in the right
places.} Counting pre-tokens is a map--reduce: each worker counts pre-tokens in a slice of the
file, and the counters are summed. The only subtlety is the slice boundaries. A cut in the middle
of a word would create two fake pre-tokens. Special tokens already define hard boundaries (no merge
may cross a document separator), so cut \emph{at} separators: pick evenly spaced byte
offsets, then move each one forward to the next occurrence of the separator, reading ahead in
small blocks (\cref{fig:proj1:chunks}). Then:
\begin{itemize}
  \item each worker opens the file itself and reads only its byte range, so you send offsets to
        the workers rather than text;
  \item within a chunk, split on \emph{all} special tokens (an alternation of
        \code{regex.escape}d strings) before applying the pattern, and do not count the special
        tokens;
  \item return a \code{Counter} of pre-token \emph{byte strings}; merging $k$ counters is cheap
        compared with the regex work.
\end{itemize}

\begin{figure}[t]
\centering
\begin{tikzpicture}[x=0.68cm, y=0.6cm]
\fill[fdata!12] (0,0) rectangle (20,1);
\foreach \x in {2.3,5.1,7.7,9.2,12.6,14.1,16.9,18.4} {\fill[warnred!50] (\x,0) rectangle +(0.25,1);}
\draw[black!50] (0,0) rectangle (20,1);
\node[font=\scriptsize, anchor=east] at (-0.1,0.5) {file};
\foreach \x/\l in {5/$b_1$,10/$b_2$,15/$b_3$} {\draw[dashed, black!60] (\x,-0.3)--(\x,1.3); \node[font=\scriptsize] at (\x,1.65) {\l};}
\foreach \x/\y in {5/5.1,10/12.6,15/16.9} {\draw[-{Stealth[length=1.6mm]}, sysblue, line width=0.7pt] (\x,-0.55) -- (\y,-0.55); \draw[sysblue, line width=0.9pt] (\y,-0.9)--(\y,0);}
\node[font=\scriptsize, sysblue, anchor=north] at (11.3,-0.9) {snap each guess forward to the next separator};
\draw[decorate, decoration={brace, amplitude=3pt, mirror}] (0,-1.8)--(5.1,-1.8) node[midway,below=3pt,font=\scriptsize]{worker 0};
\draw[decorate, decoration={brace, amplitude=3pt, mirror}] (5.1,-1.8)--(12.6,-1.8) node[midway,below=3pt,font=\scriptsize]{worker 1};
\draw[decorate, decoration={brace, amplitude=3pt, mirror}] (12.6,-1.8)--(16.9,-1.8) node[midway,below=3pt,font=\scriptsize]{worker 2};
\draw[decorate, decoration={brace, amplitude=3pt, mirror}] (16.9,-1.8)--(20,-1.8) node[midway,below=3pt,font=\scriptsize]{worker 3};
\node[font=\scriptsize, warnred, anchor=south west] at (0,1.1) {red: \texttt{<|endoftext|>}};
\end{tikzpicture}
\caption{Parallel pre-tokenisation. Evenly spaced guesses $b_i$ are moved to the next document
separator, so no pre-token or merge can straddle two workers. Chunks have unequal sizes; use more
chunks than processes to balance the load.}
\label{fig:proj1:chunks}
\end{figure}

\paragraph{The merge loop is sequential, so make each iteration incremental.} After
pre-tokenisation the corpus is a dictionary from pre-token (a tuple of byte strings) to count,
typically with millions of distinct entries. Recounting all pairs after every merge costs
$O(\text{corpus})$ per merge. Instead, maintain an index from each pair to its count and to the
pre-tokens that contain it. A merge then touches only the pre-tokens that contain the merged pair,
and adjusts counts for their neighbouring pairs. Finding the maximum can be a linear scan over the
pair-count table (often fast enough) or a heap with lazy deletion. Ties must be broken
deterministically; this project prefers the lexicographically greater pair, compared as
\code{(bytes, bytes)}. Profile before optimising (\code{cProfile}, or \code{py-spy} for a sampling
profile of a running process). On large corpora the regex pass usually dominates until you
parallelise it, and the merge loop dominates afterwards.

\begin{napkinmath}[Is your trainer fast enough?]
Measure pre-tokenisation throughput on a 100\,MB slice with one process, in MB/s. A
single-threaded regex pass typically manages a few to a few tens of MB/s in pure Python. The
time for the full file is $\text{size}/(\text{throughput}\times P)$ with $P$ workers, plus the
merge loop, whose cost depends on the number of \emph{distinct} pre-tokens (not the file size)
times the number of merges times the average number of pre-tokens each merge touches.
Extrapolate both terms from small runs before you launch a multi-hour job, and decide which one
to optimise.
\end{napkinmath}

\subsection{Encoding: special tokens, streaming and serialisation}\label{sec:proj1:stream}

\paragraph{Special tokens when encoding.} Split the input on special tokens first, emit their
ids directly, and BPE-encode only the text in between. When special tokens overlap (one is a
prefix or substring of another, such as a doubled end-of-text marker registered as its own
token), the longer one must win, so build the alternation longest first.

\paragraph{Streaming with \code{encode\_iterable}.} A generator that reads an iterable of strings
(a file handle yields lines) and yields ids lazily lets you tokenise a file that does not fit in
memory, in constant memory. It is correct only if its output equals \code{encode} on the
concatenated text, which holds as long as you cut the stream only where no pre-token or special
token can straddle the cut. Lines are a convenient unit, but think about what the pattern does
with the newline at the end of one line and the start of the next, and about special tokens split
across reads. Keep the per-chunk buffer bounded.

\paragraph{Throughput.} Real corpora repeat pre-tokens heavily, so caching pre-token $\to$ ids
(a plain dict or an LRU cache) typically speeds encoding up several-fold. Measure bytes per second
on a sample and extrapolate to your corpus, and to a terabyte-scale corpus to see why
production tokenizers are written in Rust. Use multiprocessing across documents for the full
datasets.

\paragraph{Serialising vocabularies and merges.} Token strings are \code{bytes}, which are not
JSON-safe. Choose a lossless format (pickle, base64 strings in JSON, or GPT-2's printable
byte-to-Unicode mapping) and write a round-trip test. Token-id arrays go to disk as \code{np.uint16}
(\cref{sec:optim:data}); assert that the vocabulary fits that dtype, and memory-map the arrays
during training.

\subsection{Conventions}\label{sec:proj1:conventions}

The interface fixes a few arbitrary choices so the tests can compare numbers:
\begin{itemize}
  \item \textbf{RoPE pairing}: rotate-half, as in \code{llmsys.layers.apply\_rope} and most
        released checkpoints. The adjacent-pairs layout $(2k, 2k{+}1)$ is equally valid
        (\cref{sec:modern:rope} shows the two are a fixed permutation of each other), but the test
        checks rotate-half. RoPE takes explicit integer positions.
  \item \textbf{Mask polarity}: \code{attention} takes \code{True} = allowed, like PyTorch's
        boolean \code{attn\_mask}. \code{llmsys.layers.causal\_mask} uses the opposite polarity
        internally (\code{True} = forbidden). Mixing the two silently inverts the causal mask, which
        shows up as a suspiciously low training loss.
  \item \textbf{Weight orientation}: store $W\in\R^{d_\text{out}\times d_\text{in}}$ and apply it
        as $xW^\top$, or with \code{einsum} (\cref{sec:archdesign:einsum}).
  \item \textbf{Initialisation}: truncated normals scaled by fan-in and fan-out, with unit
        embeddings and unit RMSNorm gains, are a good default (\cref{sec:archdesign:init}).
\end{itemize}

% =====================================================================
\section{Sanity checks before you train}\label{sec:proj1:milestones}

Do not move on until each check passes.

\begin{enumerate}
  \item \textbf{Tokenizer.} Train by hand on a five-word corpus and check the merges against your
        pen-and-paper trace. All M1 acceptance tests pass. Round trips are exact on a full
        TinyStories validation file.
  \item \textbf{Components.} Every M2/M3 test passes. Annotate tensor shapes
        (\code{jaxtyping} or assertions) at module boundaries.
  \item \textbf{Model at initialisation.} The loss on random data is close to $\ln V$
        (\cref{sec:optim:loss}): about 9.21 for a 10K vocabulary. The parameter count matches your
        hand count.
  \item \textbf{Overfit a single batch} of real data to near-zero loss.
  \item \textbf{Throughput.} Measure tokens/s and compare with $6N$ FLOPs per token and your
        device's peak (\cref{sec:optim:account}). If you are far off, profile: the usual culprits
        are the data loader, too much evaluation, a host sync every step, or unbatched attention.
  \item \textbf{Resume.} Kill a run, resume from its checkpoint, and compare the loss curves.
        Lab~12 automates this check for the book's loop.
\end{enumerate}

% =====================================================================
\section{Experiments}\label{sec:proj1:experiments}

Use this base configuration unless an experiment says otherwise: vocabulary 10,000, context 256,
$d=512$, 4 layers, 8 heads, $\dff=1344$ (the multiple of 64 nearest to $\tfrac83 d$), RoPE
$\theta=10^4$, AdamW with $\beta=(0.9,0.95)$, warmup + cosine to $\eta_{\max}/10$, clip 1.0. That
is about 12.5M non-embedding parameters. Training budgets are stated in tokens: $2.5\times10^8$
on a GPU, $3\times10^7$ on a laptop. \Cref{fig:proj1:plan} shows a sensible order.

\begin{figure}[t]
\centering
\begin{tikzpicture}[
  ph/.style={draw, rounded corners=2pt, minimum height=0.8cm, align=center, font=\scriptsize, line width=0.6pt},
  arr/.style={-{Stealth[length=1.8mm]}, line width=0.6pt}]
\node[ph, fill=fdata!18, draw=keygreen, text width=2.3cm] (a) at (0,0) {E1 tokenizers;\\debug config\\(1--2 min)};
\node[ph, fill=fcompute!18, draw=sysblue, text width=2.6cm] (b) at (3.3,0) {E2 coarse LR sweep\\$\times\sqrt{10}$ grid,\\short budget};
\node[ph, fill=fcompute!18, draw=sysblue, text width=2.6cm] (c) at (6.8,0) {E2 fine sweep\\full budget,\\$\ge$1 divergent run};
\node[ph, fill=fmemory!20, draw=napkin, text width=2.6cm] (d) at (10.3,0) {E3 batch sweep\\(re-tune LR per $B$,\\equal tokens)};
\node[ph, fill=fcomm!16, draw=labpurple, text width=2.6cm] (e) at (3.3,-1.8) {E4 samples; E5\\ablations, one\\change each};
\node[ph, fill=fcomm!16, draw=labpurple, text width=2.6cm] (f) at (6.8,-1.8) {E6 OpenWebText\\with the same recipe};
\node[ph, fill=warnred!12, draw=warnred, text width=2.6cm] (g) at (10.3,-1.8) {E7 60-minute\\speed run};
\draw[arr] (a)--(b); \draw[arr] (b)--(c); \draw[arr] (c)--(d);
\draw[arr] (c.south) -- ++(0,-0.35) -| (e.north); \draw[arr] (e)--(f); \draw[arr] (f)--(g);
\end{tikzpicture}
\caption{A sensible order for the experiments. Every arrow carries forward the best settings
found so far; every box writes a run directory with config, JSONL metrics (step and wall clock)
and a line in your experiment log.}
\label{fig:proj1:plan}
\end{figure}

\begin{description}[leftmargin=1.2em, style=nextline]
  \item[E1: tokenizers.] Train a 10K-vocabulary tokenizer on TinyStories and a 32K one on your
    OpenWebText sample. Report time, peak memory and the profile breakdown, and inspect the longest
    tokens. Measure compression (bytes per token) on held-out documents for each tokenizer
    \emph{and} across domains. Measure encoding throughput and extrapolate to your full corpora.
    Encode the train and validation splits to \code{uint16} files.
  \item[E2: learning rate.] Sweep the peak learning rate on a log grid, first with short runs and
    then with full-budget runs around the best value. Include runs that diverge, and relate the
    divergence point to the optimum (\cref{fig:optim:lrsweep}). Plot validation loss against
    both steps and wall-clock time.
  \item[E3: batch size.] Vary $B$ from 1 to the memory limit at equal \emph{tokens}, re-tuning
    the learning rate for each $B$ with at least three points. Discuss loss per token and loss per
    second (\cref{sec:optim:debug}, \cref{ch:scalinglaws}).
  \item[E4: decoding.] Generate 256 tokens from your best TinyStories model at several
    temperatures and top-$p$ values. Characterise the failure modes at each extreme.
  \item[E5: ablations.] Train four variants with one change each: no RMSNorm anywhere; post-norm
    blocks; no position embeddings at all (NoPE); and a non-gated SiLU MLP with $\dff=4d$, chosen so
    the parameter count matches the SwiGLU baseline. Re-tune the learning rate where the change
    plausibly moves it. \Cref{ch:archdesign} explains what each ablation probes, and Lab~7 lets you
    rehearse the methodology at toy scale.
  \item[E6: web text.] Train the base recipe on OpenWebText with the same budget. Explain the loss
    gap to TinyStories (\cref{sec:proj1:owt}) and compare samples.
  \item[E7: speed run.] See \cref{sec:proj1:speedrun}.
\end{description}

\paragraph{Logging and the report.} Every run gets a directory with its config, metadata (git
revision, hardware, versions) and metrics with step and wall time (\cref{sec:optim:explog}). Keep a
running experiment log: hypothesis, change, command, outcome, conclusion. Plots label axes with
units (nats/token), state the token budget, and use consistent colours for the same setting across
figures. \code{budget.py} keeps the arithmetic honest: steps $=\lceil
\text{tokens}/(BT)\rceil$, and the schedule must end at that step for every run
(\cref{sec:optim:sched}).

% =====================================================================
\section{Web text and the speed-run challenge}\label{sec:proj1:owt}

\paragraph{Interpreting the gap.} The same model and token budget give a much higher loss on
OpenWebText than on TinyStories. Loss in nats per token is not comparable across datasets (or
across tokenizers): it measures the model's cross-entropy, which includes the data's inherent
entropy, and web text is far more diverse than synthetic children's stories. A useful framing is
$\text{loss} = \text{irreducible entropy} + \text{model and optimisation gap}$
(\cref{ch:scalinglaws}); OpenWebText raises both terms. \Cref{ch:data} discusses what makes web
data hard and how filtering changes it.

\paragraph{The speed run.}\label{sec:proj1:speedrun} Rules: minimise the validation loss on your
held-out OpenWebText file, measured on its first two million tokens, within 60 minutes of
wall-clock time on one GPU, training only on your OpenWebText training split. Everything else is
open. It is an optimisation problem over the whole system, and the book suggests the axes to
explore. We give directions, not a recipe:
\begin{itemize}
  \item \textbf{Throughput first.} Every percent of MFU is a percent more tokens: bf16 autocast,
        \code{torch.compile}, fused attention (SDPA/FlashAttention), a prefetching sampler, no
        per-step host syncs, and evaluation only at the end (\cref{ch:perf,sec:optim:data}).
  \item \textbf{Shape for the hardware.} Dimensions that are multiples of 64 or 128, a vocabulary
        padded to such a multiple, the largest batch that stays efficient, and a model size chosen for
        the budget rather than as large as possible (\cref{ch:scalinglaws,sec:archdesign:hparams}).
  \item \textbf{Optimisation.} A schedule that ends exactly at the deadline (WSD or cosine),
        warmup, $\beta_2$, weight decay on the right parameter groups, and optimizers such as Muon
        for hidden matrices (\cref{sec:optim:beyond}).
  \item \textbf{Architecture tweaks with good evidence at small scale}: weight tying with an
        adjusted embedding initialisation, QK-norm, z-loss and other stability tricks that allow
        larger learning rates (\cref{sec:archdesign:stability}). Treat each as a hypothesis and
        ablate it.
  \item \textbf{Respect the clock.} Enforce the limit in code, stopping with time left for the
        final evaluation (\ident{code/projects/p1_lm/wallclock.py}), and plot validation loss
        against wall-clock time.
\end{itemize}
Tricks that win a one-hour contest do not necessarily help at larger scale. Project~3
(\cref{ch:proj3}) studies how to tell the difference.

% =====================================================================
\section{Hardware substitutions}\label{sec:proj1:outside}

\begin{table}[h]
\centering\small
\begin{tabularx}{\linewidth}{@{}lX@{}}
\toprule
\textbf{Hardware} & \textbf{How to adapt}\\
\midrule
CPU only & Use the laptop budget ($3\times10^7$ tokens). Set \code{torch.set\_num\_threads} to your
           physical core count, use \code{torch.compile(model)}, and train BPE on the validation
           split first.\\
Apple silicon (MPS) & As above with \code{device="mps"}. Prefer full fp32 matmuls on MPS, and if
           Inductor compilation is unavailable, use the \code{aot\_eager} backend. Tokenizer work is
           CPU-bound either way.\\
Consumer GPU (12--24\,GB) & The full budget fits. Use bf16 autocast if supported, and reduce the
           batch with gradient accumulation if needed (\cref{sec:pretrain:efficiency}).\\
Cloud H100/A100 & Rent by the hour, keep data on local NVMe, checkpoint to persistent storage,
           and script everything so a preempted instance can resume (\cref{sec:optim:ckpt}).\\
\bottomrule
\end{tabularx}
\caption{Hardware substitutions for Project~1.}
\end{table}

If you scale down, scale \emph{consistently}: fewer tokens and, if necessary, a smaller model,
with the same experiment design and logging. The qualitative findings (where the learning rate
diverges, which ablations hurt) usually survive the reduction; the absolute losses do not.

% =====================================================================
\section{Summary}
\begin{itemize}
  \item Project~1 builds a tokenizer, a modern Transformer LM, an optimizer and a training system
        from first principles, then studies them experimentally.
  \item A thin interface file connects your code to the acceptance tests, which compare against
        PyTorch, the \code{llmsys} reference and behavioural properties.
  \item BPE at scale is a parallel map--reduce over chunks cut at special tokens, followed by a
        sequential, incrementally maintained merge loop. Streaming encoding may cut only where no
        token can straddle.
  \item Fix the token budget, derive steps from it, end the schedule there, and log step and
        wall-clock time for every run.
  \item For the speed run, measure and optimise throughput before you tune the model, and enforce
        the time limit in code.
\end{itemize}

\section{Exercises}

\begin{exercise}[Pre-token boundaries]
Apply the GPT-2 pattern by hand to the string \verb*|x  = 10.5;  y| and list the pre-tokens.
Which alternative of the pattern produced each one? Then check with the \code{regex} module.
\end{exercise}

\begin{exercise}[Chunk balance]
A 12\,GB file has documents whose lengths are roughly exponential with mean 4\,KB. You cut it
into $k$ chunks at document separators and process them with 8 workers. Estimate the largest
chunk-size deviation for $k=8$ and $k=64$, and explain why more chunks than workers help.
\end{exercise}

\begin{exercise}[Accounting on another model]
Using \cref{sec:optim:account}, compute the parameter count and forward matmul FLOPs per token
for a pre-norm SwiGLU model with $V=32{,}000$, $L=16$, $d=1024$, 16 heads, $\dff=2816$ and untied
embeddings. What fraction of the forward FLOPs is the attention-score computation ($QK^\top$ and
$PV$) at $T=2048$ and at $T=32{,}768$?
\end{exercise}

\begin{exercise}[Budgeted sweep design]
You have 6 GPU-hours, a 20-minute full-budget run, and 4-minute short runs that preserve the
learning-rate ranking. Design an E2 + E3 study that fits the budget, includes a divergent run,
and re-tunes the learning rate per batch size.
\end{exercise}
